%% =====================================================================
\section{Related Work}
\label{sec:related}

\paragraph{The cube group and its diameter.}
The \Cu{3} group and its generalisations are classical objects of group theory~\cite{singmaster1981notes,joyner2008adventures}.
Upper bounds on the \Cu{3}'s God's number fell from Thistlethwaite's 52 face turns to 26 by
Kunkle and Cooperman's disk-based search~\cite{kunkle2007twentysix}, and Rokicki, Kociemba, Davidson and
Dethridge settled it at 20 face turns by solving all $2.2\cdot10^9$ cosets of a subgroup on donated
computer time~\cite{rokicki2014diameter}; in the quarter-turn metric it is 26~\cite{rokicki2014qtm}. For larger
cubes the diameter is unknown. Demaine et al.~\cite{demaine2011algorithms} prove that it is $\Theta(N^2/\log N)$
for both the \Cu{N} and the \Sq{N}: the lower bound is a counting argument (there are $e^{\Theta(N^2)}$ states
but only $\Theta(N)$ moves, so some state needs $\Omega(N^2/\log N)$ of them), and the matching upper bound
solves $\Theta(\log N)$ clusters with each shared move sequence. Ellenberg~\cite{ellenberg2011diameter}
explains why the counting bound is the natural guess, and the CayleyPy authors compute growth functions of
many Cayley graphs to conjecture diameters~\cite{chervov2025growth}. Demaine, Eisenstat and Rudoy show that
deciding whether a \Cu{N} state can be solved in $k$ moves is NP-complete~\cite{demaine2018npcomplete}. We use
the counting argument quantitatively (\S\ref{sec:bound}): with exact state counts and the exact growth of
canonical quarter-turn words it gives a concrete bound for every $N$ that we can put next to every method.

\paragraph{Optimal and near-optimal \Cu{3} solvers.}
Thistlethwaite's algorithm and Kociemba's two-phase algorithm~\cite{kociemba2phase} solve the \Cu{3} by
descending through subgroups; two-phase returns about 20 face turns in milliseconds, and min2phase is a fast
engineered implementation~\cite{chen2phase}. Korf found the first optimal solutions of random states with
IDA* and pattern databases~\cite{korf1997optimal}, building on Culberson and Schaeffer's pattern
databases~\cite{culberson1998pattern}. Rokicki's \texttt{nxopt}~\cite{rokicki2018nxopt} refines this with coset
pruning tables of up to terabytes, 16-way symmetry and inverse lookups, in the half-turn and the quarter-turn
metric; Kociemba's Cube Explorer also contains a quarter-turn optimal solver~\cite{kociemba2qtm}. \Mone is our
reproduction of Korf's solver, and \Mseven uses \texttt{nxopt} for the 3$\times$3$\times$3 stage of odd cubes.

\paragraph{Large-cube solvers.}
Human and computer solvers for big cubes use \emph{reduction}: solve the centers, pair the edges, and finish
as a \Cu{3}~\cite{walton2017nxnxn}. Korf's macro-operators~\cite{korf1985macro} formalise the move sequences
that change one part of the state and restore the rest; commutators $[A,B]=ABA'B'$ are the standard way to
build them, and conjugates $sCs'$ move their effect anywhere. Demaine et al.'s parallel algorithm batches
clusters that are in the \emph{same} configuration so that one sequence of slice moves solves all of them;
at practical sizes a center cluster has about $3\cdot10^{15}$ configurations, so exact matches are rare, and
\Mthree and \FiveC batch 3-cycles instead. The dwalton76 solver~\cite{walton2017nxnxn} is a mature open
reduction solver with large lookup tables for $N\le7$ and a general path above; we run it as a baseline.

\paragraph{Kaggle Santa 2023 and KMCoders.}
The Santa 2023 competition~\cite{kaggle2023santa} asked for short solutions of 398 permutation puzzles: cubes
from \Cu{2} to \Cu{33}, ``wreaths'' and ``globes'', each given as a start state, a goal state and a set of
allowed moves, scored by the total number of moves (\S\ref{sec:data}). KMCoders published their solution,
their final submission and their Rust code~\cite{kmcoders2024santa}. For cubes they solve the special parts
(corners, and for odd $N$ the central \Cu{3}) optimally, make every remaining cluster an even permutation and
roughly align it with simulated annealing over elementary moves~\cite{kirkpatrick1983annealing}, then insert
short commutators and finally 3-cycles at \emph{arbitrary} positions of the move sequence, choosing time and
sequence greedily or by beam search; their code also applies bulk commutators of slice sets chosen by
annealing. Their reported lengths are the best of many runs. We reproduce the pipeline from the write-up
(\Msix), run their program unchanged as a baseline, and build \Mseven around it.

\paragraph{Learned solvers.}
DeepCubeA~\cite{agostinelli2019deepcubea} learns a cost-to-go function by approximate value iteration and
searches with batch weighted A*. EfficientCube~\cite{takano2023selfsupervision} trains directly on random
scrambles from the goal. Q* search~\cite{agostinelli2024qstar} scores all children of a node in one network
pass and handles large action spaces. CayleyPy~\cite{chervov2025cayleypy} trains a residual MLP to predict the
length of a random walk from the solved state and searches with very wide beams; it reports over 98\%
optimal \Cu{3} solutions and the first learned \Cu{4} and \Cu{5} solvers. All of these models are trained per
puzzle size on the full sticker array. \Mfour is our reproduction of CayleyPy; \FiveA instead
decompose the puzzle into clusters whose networks are shared by all sizes.

\paragraph{Search techniques we build on.}
Beam search~\cite{lowerre1976harpy} keeps the $W$ best nodes per depth; \Mfour and \FiveA use it, as
does KMCoders' final stage. Simulated annealing~\cite{kirkpatrick1983annealing} is KMCoders' second stage and
\Msix's. Large-neighbourhood search~\cite{shaw1998lns}, which destroys and rebuilds part of a solution, is
what we tried, without success, on top of KMCoders' greedy stage (\S\ref{sec:m7}). \Mseven belongs to the
line of work on randomisation and portfolios: run times and solution qualities of search algorithms are
heavy-tailed over random seeds, so restarts and portfolios of independent runs pay off~\cite{gomes1998boosting,
huberman1997portfolio,gomes2001portfolios}. The cube's symmetries are usually used to \emph{reduce} search,
by storing one representative per symmetry class in pattern databases and coset tables~\cite{rokicki2014diameter,
rokicki2018nxopt}; \Mseven uses them the other way round, as a source of diverse but equivalent inputs to a
solver whose own randomness is fixed.

\paragraph{Size generalisation and other uses of the cube.}
Networks that read a problem as a set or sequence of parts, such as pointer
networks~\cite{vinyals2015pointer} and set transformers~\cite{lee2019settransformer}, can be applied to larger
instances than they were trained on; \FiveA applies this idea to cubie clusters. The cube also serves as a
benchmark for robotic manipulation~\cite{openai2019dactyl} and for spatial reasoning in (multimodal) language
models~\cite{anand2025cubebench,gao2025cubebench}; those works use small cubes and do not target solution
length at scale.
